\section*{Chapter \ref{ch:rc:mortgage-modelling} --- Mortgage Modelling}

\begin{solution}{exo:rc:mortgage-modelling:1}
The mortgage rate is $3.80+1.75 = 5.55\%$; the incentive $6.00-5.55 = 0.45$ percentage point,
just below the S-curve's centre.
\end{solution}

\begin{solution}{exo:rc:mortgage-modelling:2}
$181-146 = 34$ basis points (35 before rounding the two spreads): the part of the pool's
yield over the curve that pays for the borrowers' prepayment options once rates are
volatile, measured as a spread.
\end{solution}

\begin{solution}{exo:rc:mortgage-modelling:3}
Higher rates slow refinancing, so the balance on which interest is paid survives longer;
the extra interest more than offsets the higher discounting.
\end{solution}

\begin{solution}{exo:rc:mortgage-modelling:4}
Prepayment depends on the path (\omterm{def:rc:mortgage-modelling:scurve}{burnout}, the remaining balance), not only on the current
rate: two paths reaching the same node carry different pools. A recombining tree forgets
the path; simulation keeps it.
\end{solution}

\begin{solution}{exo:rc:mortgage-modelling:5}
A fall of 100 basis points gains about 4.54, a rise loses about 5.50: the losses exceed the
gains, the asymmetry of negative convexity.
\end{solution}

\begin{solution}{exo:rc:mortgage-modelling:6}
Without \omterm{def:rc:mortgage-modelling:scurve}{burnout} borrowers keep refinancing at the rate the S-curve says, so the premium
pool (coupon above the market's) returns principal at par faster, cutting off the
above-market coupon earlier: its value falls, to 99.20.
\end{solution}

\begin{solution}{exo:rc:mortgage-modelling:7}
DV01 before $4.99\times100/100\times10\text{ billion}\times10^{-4} = \text{USD}~4.99$ million per basis
point, after $6.17\times94.50/100\times\ldots = 5.83$ million; the extra USD~844\,025 per basis
point divided by the ten-year swap's USD~796\,282 per billion gives USD~1.06 billion of
payer swaps.
\end{solution}

\begin{solution}{exo:rc:mortgage-modelling:8}
Duration hedging removes first-order risk only: the book is short convexity (and short
volatility through the borrowers' options), so it loses on large moves in either direction
between re-hedges, pays bid--offer on every re-hedge, and carries prepayment-model and
mortgage-spread risk that no rate hedge covers.
\end{solution}

\begin{solution}{pb:rc:mortgage-modelling:1}
\textbf{1.} 146, 181 and 34 basis points.
\textbf{2.} 4.99 years and $-123$.
\textbf{3.} USD~4.99 million per basis point.
\textbf{4.} Interest 25.26, principal 74.74 per 100.
\textbf{5.} The pool: it is short the borrowers' options to prepay, sold implicitly by
lending to them at a fixed rate with a free prepayment right.
\textbf{6.} 94.50.
\textbf{7.} 6.17 years; USD~5.83 million per basis point.
\textbf{8.} Higher rates move the pool out of the refinancing S-curve: prepayments slow and
the principal comes back later.
\textbf{9.} USD~1.06 billion of ten-year payer swaps.
\textbf{10.} Their payer swaps push swap rates up further, which lengthens the pools again,
which requires more hedging: a feedback the Federal Reserve study found in the data.
\textbf{11.} IO 24.00 at $-25$ and 26.61 at $+25$; PO 77.21 and 72.10.
\textbf{12.} The IO, whose value rises with rates, offsets part of the PO's (and the pool's)
negative convexity.
\textbf{13.} Deposits that reprice slowly are worth more when rates rise; an IO behaves the
same way and can hedge the bank's balance sheet (chapter~24).
\textbf{14.} Value: the principal arrives later than priced, at the same par amount.
\textbf{15.} The IO: its value depends almost entirely on how long the balance survives.
\textbf{16.} Swaptions would buy back convexity and cost option premium (the pool's \omterm{def:rc:mortgage-modelling:optioncost}{option
cost} is the benchmark); many holders mix both.
\textbf{17.} Duration extends gradually as rates rise and re-hedging flows spread over weeks;
volatility rises and feeds back, and the Fed study found the effects persist for months.
\textbf{18.} A holder that does not hedge removes part of the feedback; one that stops
buying (runs off its holdings) shifts the hedging to private holders and can raise it.
\textbf{19.} Named result: \emph{the convexity event}: on USD~10 billion of the pool a 100-basis-point
rise lifts DV01 from USD~4.99 million to 5.83 million per basis point, and re-hedging takes
USD~1.06 billion of ten-year payer swaps.
\textbf{20.} Because borrowers stop refinancing when rates rise, exactly when the holder
would like its money back.
\end{solution}

\begin{solution}{iq:rc:mortgage-modelling:1}
Borrowers can prepay at par whenever rates fall: the security's price gains are capped as
rates fall (it shortens) and its losses grow as rates rise (it lengthens). It is a bond
minus a call option held by the borrowers.
\iqlookfor{the embedded prepayment option and asymmetric price moves.}
\end{solution}

\begin{solution}{iq:rc:mortgage-modelling:2}
Turnover (home sales, defaults, curtailments; age-driven, little rate dependence);
refinancing (an S-curve in the incentive, WAC minus mortgage rate, damped by \omterm{def:rc:mortgage-modelling:scurve}{burnout}
and pool characteristics); the mortgage rate from market rates through the
\omterm{def:rc:mortgage-modelling:cc}{primary--secondary spread}. Refinancing and the mortgage rate depend on rates.
\iqlookfor{the components and which are rate-driven.}
\end{solution}

\begin{solution}{iq:rc:mortgage-modelling:3}
Simulate rate paths from a calibrated term-structure model; run the \omterm{def:rc:mortgage-modelling:model}{prepayment model} and
the pool's cash flows on each path; find the constant spread over the path's short rate
that makes the average discounted cash flow equal the price. Monte Carlo because
prepayment is path-dependent.
\iqlookfor{paths, cash flows, spread root-finding, and path dependence.}
\end{solution}

\begin{solution}{iq:rc:mortgage-modelling:4}
Their pools lengthen; to keep duration they pay fixed in swaps or sell Treasuries, adding
to the sell-off; implied volatility rises as hedgers buy options; the move can overshoot
before the flows exhaust.
\iqlookfor{hedging flows in the direction of the move.}
\end{solution}

\begin{solution}{iq:rc:mortgage-modelling:5}
Run several \omterm{def:rc:mortgage-modelling:model}{prepayment models} or perturbed parameters (speed multipliers, S-curve shifts,
\omterm{def:rc:mortgage-modelling:scurve}{burnout}), compute the option-adjusted spread and risk under each, report prepayment
durations, and reserve for the dispersion of values; back-test the model against realised
speeds by cohort.
\iqlookfor{scenario dispersion and back-testing.}
\end{solution}

\begin{solution}{iq:rc:mortgage-modelling:6}
Random numbers not common across the bumped runs (a new seed each time), too few paths,
path sets that change size, the curve or volatility inputs moving, and the finite-difference
bump interacting with a steep S-curve; fix seeds and paths, use antithetics, and check the
bump size.
\iqlookfor{common random numbers and noise diagnosis.}
\end{solution}
